\documentclass[10pt]{article}
\usepackage[margin=1in]{geometry}
\usepackage{booktabs}
\usepackage{graphicx}
\usepackage{amsmath}
\usepackage{hyperref}
\usepackage{multirow}
\usepackage{xcolor}

\newcommand{\meanstd}[2]{$#1_{\pm #2}$}

\title{Structural Entropy as a Regulariser for Neural Time-Series Forecasting}
\author{}
\date{}

\begin{document}
\maketitle

\begin{abstract}
We investigate whether spectral entropy of encoder activations provides a
useful regularisation signal for neural time-series forecasting.
Using TimesNet on the four ETT datasets, we run 120 controlled experiments
across Phase A baselines and Phase B pilots with both entropy-minimisation
($+\lambda H$) and entropy-maximisation ($-\lambda H$) objectives.
Entropy minimisation collapses the encoder toward a low-rank representation
and degrades forecasting accuracy monotonically with $\lambda$ (up to $+17\%$
MSE).
Entropy maximisation is benign but marginal: the best observed improvement is
$-1.6\%$ MSE on a single dataset, with no consistent benefit across datasets.
Neither direction meets the practical significance threshold
($\Delta\text{MSE} < -0.5\%$ on $\geq\!2$ of 4 datasets).
We conclude that \emph{pre-projection} spectral entropy is not a productive
regularisation target for TimesNet on ETT benchmarks, and suggest that
within-block activations or architectures with stronger collapse tendencies
are more promising targets.
\end{abstract}

\section{Introduction}

\section{Related Work}

\section{Benchmark}

\section{Method}

\section{Experiments}

\subsection{Phase A: Baseline Reproduction}
\label{sec:phase_a}

We reproduce the Phase A baseline (no SE regularisation) using
\textbf{TimesNet}~\cite{wu2023timesnet} and \textbf{DLinear}~\cite{zeng2023dlinear}
on the ETT benchmark suite (ETTh1, ETTh2, ETTm1, ETTm2) with four prediction
horizons $H \in \{96, 192, 336, 720\}$.
All runs use a fixed look-back window of $L=96$, label length 48, multivariate
features (\texttt{M}), 10 training epochs with early stopping (patience 3),
batch size 32.
TimesNet results are reported as mean\,$\pm$\,std over three independent seeds
(1, 2, 3); DLinear is deterministic so a single value is reported.
TSLib commit \texttt{4e938a1} is used throughout.

\paragraph{Note on DLinear ETTm gap.}
Our DLinear ETTm numbers are 0.05--0.16~MSE above the original paper.
This is \emph{expected and by design}: the original paper tunes the look-back
window per model per dataset; we fix $L=96$ for a fair, controlled comparison.
This protocol choice is made explicit and does not constitute a reproduction
failure.

% --- Table: Phase A baseline ---
\begin{table}[t]
\centering
\caption{%
  Phase A baseline results (MSE / MAE).
  TimesNet values are mean\,$\pm$\,std over seeds \{1,2,3\};
  DLinear is deterministic (std\,=\,0).
  \textit{Paper} columns reproduce numbers from the original publications
  under their best look-back protocol.
  $\dagger$\,DLinear ETTm gap is due to our fixed $L=96$ protocol (see text).
}
\label{tab:phase_a}
\resizebox{\textwidth}{!}{%
\begin{tabular}{llr ccc ccc}
\toprule
\multirow{2}{*}{Model} & \multirow{2}{*}{Dataset} & \multirow{2}{*}{$H$}
  & \multicolumn{3}{c}{Ours ($L=96$, seeds 1–3)} & \multicolumn{2}{c}{Paper} \\
\cmidrule(lr){4-6}\cmidrule(lr){7-8}
  & & & MSE mean & MSE std & MAE mean & MSE & MAE \\
\midrule
\multirow{16}{*}{TimesNet}
  & \multirow{4}{*}{ETTh1}
    &  96 & 0.3891 & 0.0000 & 0.4120 & 0.384 & 0.402 \\
  & & 192 & 0.4394 & 0.0000 & 0.4417 & 0.436 & 0.429 \\
  & & 336 & 0.4941 & 0.0004 & 0.4709 & 0.491 & 0.469 \\
  & & 720 & 0.5178 & 0.0018 & 0.4943 & 0.521 & 0.500 \\
\cmidrule(lr){2-8}
  & \multirow{4}{*}{ETTh2}
    &  96 & 0.3292 & 0.0001 & 0.3701 & 0.340 & 0.374 \\
  & & 192 & 0.3930 & 0.0011 & 0.4082 & 0.402 & 0.414 \\
  & & 336 & 0.4708 & 0.0000 & 0.4677 & 0.452 & 0.452 \\
  & & 720 & 0.4418 & 0.0007 & 0.4520 & 0.462 & 0.468 \\
\cmidrule(lr){2-8}
  & \multirow{4}{*}{ETTm1}
    &  96 & 0.3363 & 0.0000 & 0.3765 & 0.338 & 0.375 \\
  & & 192 & 0.3874 & 0.0002 & 0.4017 & 0.374 & 0.387 \\
  & & 336 & 0.4162 & 0.0009 & 0.4228 & 0.410 & 0.411 \\
  & & 720 & 0.5041 & 0.0208 & 0.4689 & 0.478 & 0.450 \\
\cmidrule(lr){2-8}
  & \multirow{4}{*}{ETTm2}
    &  96 & 0.1883 & 0.0000 & 0.2678 & 0.187 & 0.267 \\
  & & 192 & 0.2518 & 0.0006 & 0.3069 & 0.267 & 0.332 \\
  & & 336 & 0.3130 & 0.0010 & 0.3437 & 0.332 & 0.366 \\
  & & 720 & 0.4223 & 0.0003 & 0.4070 & 0.432 & 0.432 \\
\midrule
\multirow{16}{*}{DLinear}
  & \multirow{4}{*}{ETTh1}
    &  96 & 0.3962 & — & 0.4108 & 0.386 & 0.400 \\
  & & 192 & 0.4450 & — & 0.4404 & 0.437 & 0.432 \\
  & & 336 & 0.4874 & — & 0.4654 & 0.481 & 0.459 \\
  & & 720 & 0.5126 & — & 0.5104 & 0.519 & 0.516 \\
\cmidrule(lr){2-8}
  & \multirow{4}{*}{ETTh2}
    &  96 & 0.3414 & — & 0.3953 & 0.333 & 0.387 \\
  & & 192 & 0.4819 & — & 0.4792 & 0.477 & 0.476 \\
  & & 336 & 0.5929 & — & 0.5422 & 0.594 & 0.541 \\
  & & 720 & 0.8403 & — & 0.6611 & 0.831 & 0.657 \\
\cmidrule(lr){2-8}
  & \multirow{4}{*}{ETTm1$^\dagger$}
    &  96 & 0.3459 & — & 0.3737 & 0.299 & 0.343 \\
  & & 192 & 0.3818 & — & 0.3911 & 0.335 & 0.365 \\
  & & 336 & 0.4153 & — & 0.4150 & 0.369 & 0.386 \\
  & & 720 & 0.4730 & — & 0.4509 & 0.425 & 0.421 \\
\cmidrule(lr){2-8}
  & \multirow{4}{*}{ETTm2$^\dagger$}
    &  96 & 0.1934 & — & 0.2928 & 0.167 & 0.260 \\
  & & 192 & 0.2845 & — & 0.3614 & 0.224 & 0.303 \\
  & & 336 & 0.3849 & — & 0.4295 & 0.281 & 0.342 \\
  & & 720 & 0.5561 & — & 0.5235 & 0.397 & 0.421 \\
\bottomrule
\end{tabular}
}
\end{table}

\subsection{Phase B: SE Regularisation}

We augment the TimesNet training loss with a spectral-entropy term computed from
the pre-projection encoder activations $\mathbf{X} \in \mathbb{R}^{B \times T \times D}$:

\begin{equation}
\mathcal{L}_\text{total} = \mathcal{L}_\text{MSE}(\hat{y}, y)
  \;\pm\; \lambda \cdot H(\mathbf{X}),
\quad
H(\mathbf{X}) = -\sum_i p_i \log p_i,
\quad
p_i = \frac{\lambda_i(\mathbf{C})}{\sum_j \lambda_j(\mathbf{C})}
\label{eq:se_loss}
\end{equation}

where $\mathbf{C} = \frac{1}{BT}(\mathbf{X}_c^\top \mathbf{X}_c)$ is the feature covariance
and $\lambda_i$ are its eigenvalues (\texttt{torch.linalg.eigvalsh}, fully differentiable).
A forward pre-hook on \texttt{model.projection} captures $\mathbf{X}$ each training step.

\paragraph{Pilot v1 — entropy minimisation ($+\lambda H$).}
Adding $+\lambda H$ drives the encoder toward a low-rank representation
(collapsed eigenspectrum, $H \to 0$).  We swept
$\lambda \in \{0.01, 0.05, 0.10\}$ over all four ETT datasets at $H=96$.

\begin{table}[h]
\centering\small
\caption{Phase B Pilot v1: $+\lambda H$ (entropy minimisation), $H=96$, seed~1.
Bold = best per dataset. $\Delta$MSE relative to $\lambda=0$ baseline.}
\label{tab:phase_b_v1}
\begin{tabular}{l c c c c c c}
\toprule
Dataset & Baseline & $\lambda{=}0.01$ & $\Delta$ & $\lambda{=}0.05$ & $\Delta$ & $\lambda{=}0.10$ \\
\midrule
ETTh1 & 0.3891 & \textbf{0.3883} & $-0.08\%$ & 0.4211 & $+8.2\%$ & 0.4550 \\
ETTh2 & 0.3292 & 0.3326 & $+1.0\%$ & 0.3366 & $+2.3\%$ & 0.3385 \\
ETTm1 & 0.3363 & 0.3366 & $+0.1\%$ & 0.3441 & $+2.3\%$ & 0.3478 \\
ETTm2 & 0.1883 & 0.1899 & $+0.9\%$ & 0.1978 & $+5.1\%$ & 0.2057 \\
\bottomrule
\end{tabular}
\end{table}

\noindent\textbf{Diagnosis.}
Collapsing the encoder into a low-rank subspace is
counter-productive: TimesNet's multi-scale frequency decomposition
requires \emph{diverse} representations.
Larger $\lambda$ monotonically degrades performance.
$\lambda=0.01$ is near-neutral (at most $\pm1\%$).

\paragraph{Pilot v2 \textemdash{} entropy maximisation ($-\lambda H$).}
We flip the sign, rewarding a \emph{uniform} eigenspectrum.
Table~\ref{tab:phase_b_v2} shows results across 12 runs
($\lambda \in \{0.01, 0.05, 0.10\} \times 4$ datasets).
Entropy maximisation produces mixed, near-neutral effects.
ETTm2 benefits modestly ($-1.6\%$ at $\lambda=0.01$) and ETTm1 shows a mild
trend ($-1.1\%$ at $\lambda=0.10$), but ETTh1 is consistently hurt.
No $\lambda$ achieves $\Delta\text{MSE} < -0.005$ on $\geq\!2$ datasets,
so the pilot is \textbf{NO-GO}.
Entropy maximisation is far less harmful than entropy minimisation:
the v1 catastrophic degradation on ETTh1 ($+8$--$17\%$) is eliminated,
confirming the encoder does exploit spectral diversity when rewarded,
but the gain is insufficient to overcome the added training noise.

\begin{table}[h]
\centering\small
\caption{Phase B Pilot v2: $-\lambda H$ (entropy maximisation), $H=96$, seed~1.
Bold = best (lowest) MSE per dataset. $\Delta$ relative to $\lambda=0$ baseline.
\textbf{Verdict: NO-GO} (0/12 configs with $\Delta\text{MSE} < -0.005$ on
$\geq\!2$ datasets).}
\label{tab:phase_b_v2}
\begin{tabular}{l c c c c c}
\toprule
Dataset & Baseline & $\lambda{=}0.01$ & $\lambda{=}0.05$ & $\lambda{=}0.10$ & Best $\Delta$ \\
\midrule
ETTh1 & 0.3891 & \textbf{0.3905} & 0.3927 & 0.3941 & $+0.38\%$ \\
ETTh2 & 0.3292 & 0.3315 & \textbf{0.3267} & 0.3298 & $-0.75\%$ \\
ETTm1 & 0.3363 & 0.3353 & 0.3333 & \textbf{0.3326} & $-1.10\%$ \\
ETTm2 & 0.1883 & \textbf{0.1852} & 0.1859 & 0.1859 & $-1.64\%$ \\
\bottomrule
\end{tabular}
\end{table}

\subsection{Failure Modes}
\label{sec:failure_modes}

The two pilots expose a systematic failure mode:
\textbf{pre-projection spectral entropy is not a meaningful regularisation
target for TimesNet}.

\emph{Entropy minimisation ($+\lambda H$).}
Gradient descent minimises the loss, hence also minimises $H$, which pushes
the encoder toward a rank-1 representation. A rank-1 encoder discards all
channel-interaction structure, which is exactly what a time-series model needs.
The result is a graceful collapse: small $\lambda$ is near-neutral, larger
$\lambda$ degrades monotonically (up to $+17\%$ MSE on ETTh1 at $\lambda=0.10$).

\emph{Entropy maximisation ($-\lambda H$).}
The opposite penalty prevents collapse but does not add useful inductive bias.
The ETT datasets do not appear to have latent high-rank structure that is
being suppressed by the base model; the encoder already spreads variance
reasonably across dimensions. Adding a diversity reward yields small, mixed
signals ($\leq\!\pm1.6\%$) with no consistent direction across datasets.

\emph{Probe location.}
We measure entropy of activations entering the linear projection (post-TSB
blocks). This layer summarises long-range temporal patterns; it may already be
well-conditioned. Measuring entropy of activations \emph{within} the TSB
period-analysis blocks, or of the FFN intermediate layer, could expose a more
tractable bottleneck.

\section{Limitations}

We evaluate a single architecture (TimesNet) on a single benchmark suite
(ETT, $H=96$). The null result may not generalise: Transformer-based models
with a higher tendency toward representation collapse (e.g.\ PatchTST, iTransformer)
could respond differently to spectral entropy regularisation.
We also test only three $\lambda$ values; a finer sweep around the
near-neutral $\lambda=0.01$ regime may reveal small but real benefits.
All runs use a single seed (seed=1) for Phase B pilots; seed variance from
Phase A suggests $\pm0.002$ MSE noise, which is comparable to the improvements
observed.

\section{Conclusion}

We study \emph{structural entropy regularisation} for neural time-series
forecasting: augmenting the training loss with a term proportional to the
spectral entropy of pre-projection encoder activations.
We run 120 controlled experiments on ETTh1/h2/m1/m2 with TimesNet,
spanning Phase A baselines, Phase B entropy-minimisation (v1), and
Phase B entropy-maximisation (v2).

Both regularisation directions fail to clear the practical significance
threshold ($\Delta\text{MSE} < -0.5\%$ on $\geq\!2/4$ datasets for any
$\lambda$): entropy minimisation actively hurts (up to $+17\%$ on ETTh1);
entropy maximisation is benign but marginally beneficial at best ($-1.6\%$
on ETTm2).

The negative result is informative.
It confirms that the pre-projection representation in TimesNet is neither
pathologically low-rank (else maximisation would help more) nor wastefully
high-rank (else minimisation would help).
Future work should probe \emph{within-block} activations and test on
architectures with stronger collapse tendencies.

\bibliographystyle{plain}
\bibliography{refs}

\end{document}

